\documentclass[11pt,a4paper]{article}

% ============================================================
% Uniform MTDF preamble -- identical across all papers
% ============================================================
\usepackage[utf8]{inputenc}
\usepackage[T1]{fontenc}
\usepackage{lmodern}
\usepackage[english]{babel}
\usepackage[a4paper,margin=2.2cm]{geometry}
\usepackage{amsmath,amssymb,amsthm}
\usepackage{graphicx}
\usepackage{float}
\usepackage{booktabs}
\usepackage{array}
\usepackage{longtable}
\usepackage{tabularx}
\usepackage{xltabular}
\usepackage{multirow}
\usepackage{caption}
\usepackage[dvipsnames,table]{xcolor}
\usepackage{mdframed}
\usepackage{parskip}
\usepackage{ragged2e}
\usepackage[hidelinks,breaklinks=true]{hyperref}
\usepackage{microtype}
\usepackage{url}
% Allow \path{} to break more aggressively inside narrow table columns:
\makeatletter
\g@addto@macro\UrlBreaks{\do\_\do\-\do\.\do\/\do\:}
\makeatother

% Colours matching the HTML theme
\definecolor{accent}{HTML}{1A4B8A}
\definecolor{callout-bg}{HTML}{FBFBFB}
\definecolor{tablehead}{HTML}{F0F0F0}

% Prevent any table from overflowing
\setlength{\tabcolsep}{4pt}
\renewcommand{\arraystretch}{1.15}

% Caption style (compact, italic, small like the HTML)
\captionsetup{font=small,labelfont=bf,labelsep=period,justification=centerfirst}

% Hyperref metadata placeholders (filled per-paper below)
\hypersetup{
  pdftitle={MTDF Companion Paper Part IV: Cosmological Validation, High Energy, and Test Programme},
  pdfauthor={Ingo Mesche},
  pdfsubject={MTDF - Mesche Time-Dilation Field Theory},
  colorlinks=false
}

% Prevent overfull hbox from long URLs/DOIs in bibliography
\sloppy
\emergencystretch=3em

% Tolerant line breaking so long identifiers (Python filenames etc.) don't
% produce large overfull hboxes in narrow table columns. These settings keep
% all content on-page while slightly relaxing right-margin strictness.
\tolerance=1200
\hbadness=2000
\hyphenpenalty=150
\exhyphenpenalty=50

% ============================================================
\title{MTDF Companion Paper Part IV: Cosmological Validation, High Energy, and Test Programme}
\author{Ingo Mesche\\ \small Independent Researcher, Malta \\ \small \texttt{imesche@outlook.com}}
\date{V75 / July 2026}

\begin{document}
\maketitle
\begin{center}\itshape Extensions, hypotheses, and test programmes\end{center}



\tableofcontents

\begin{abstract}
This companion paper (Part IV) documents the cosmological validation of MTDF through a matched CMB+BAO comparison and a full Planck plik TTTEEE + low-$\ell$ + lensing hard-falsifier, alongside high-energy and cluster-scale phenomenology. The full Planck comparison (Section 2.4, sealed v2 chains) leaves the standard parameters essentially unshifted, gives a broad coupling posterior k\textsubscript{f} = 0.685 $\pm$ 0.500 (consistency, not support), and establishes the sealed growth posture: the corrected covariant $\mu$(a) enhances late-time growth, $\sigma$\textsubscript{8} = 0.8265 $\pm$ 0.0062, about 1.9$\sigma$ (posterior units) above the same-pipeline $\Lambda$CDM baseline, so the corrected theory mildly worsens rather than eases the weak-lensing comparison; the earlier $\sigma$\textsubscript{8} = 0.790 relief was a gauge artifact and is retired. The comparison against the matched $\Lambda$CDM control is complete: $\Delta$$\chi$$^{2}$ (bounded best fits) = $-$1.29. This is a maximum-likelihood comparison, not a Bayes factor or Bayesian evidence calculation. A narrower CMB+BAO summary run in Section 2.1 ($\Delta$$\chi$$^{2}$ $\approx$ 1.25) is a separate, pre-correction bookkeeping run and is not the full-Planck total. Under the MTDF void-enhancement mapping, conditional on a deep KBC-like local underdensity, the SH0ES tension is reduced from 3.8$\sigma$ to 1.3$\sigma$ without retuning the global fit (implied local value $\approx$ 73.1 km/\allowbreak{}s/\allowbreak{}Mpc); shallower void reconstructions instead give roughly 70 to 71 km/\allowbreak{}s/\allowbreak{}Mpc, and MTDF does not predict the depth of the local void itself. We additionally present falsifiable predictions for lensing in merging clusters and stress-support contributions to hydrostatic mass bias.
\end{abstract}

\begin{mdframed}[linecolor=accent,linewidth=0.5pt,backgroundcolor=callout-bg,leftline=true,rightline=true,topline=true,bottomline=true,innerleftmargin=8pt,innerrightmargin=8pt,innertopmargin=6pt,innerbottommargin=6pt]
\textbf{Claim status.}
This paper distinguishes between:
(i) calibration anchors,
(ii) validation targets,
(iii) diagnostic consistency checks, and
(iv) speculative extensions.
Only results explicitly labelled as validation targets are used as evidence for the core MTDF claim. Diagnostic and speculative sections are included to define future tests, not to claim established confirmation.
\end{mdframed}

\section{Purpose and scope}
\label{sec:purpose}
This companion part documents two extension tracks and their tests: (i) a likelihood-based cosmological validation exercise comparing MTDF and $\Lambda$CDM on Planck 2018 TTTEEE and DESI DR1 BAO, including an explicit global-versus-local H0 interpretation when a SH0ES prior is introduced, and (ii) high-energy and cluster-scale phenomenology where stress support may affect lensing, hydrostatic mass bias, and merging systems.

The aim is to state clear, falsifiable expectations, report analyses conducted to date with full provenance, and define what would count as decisive failure modes. Where comparisons are made, the baseline choices are matched across models.

\begin{mdframed}[linecolor=accent,linewidth=0.5pt,backgroundcolor=callout-bg,leftline=true,rightline=true,topline=true,bottomline=true,innerleftmargin=8pt,innerrightmargin=8pt,innertopmargin=6pt,innerbottommargin=6pt]
Note on separation from the dashboard: the strict dashboard scorecard is intentionally conservative and does not include the full-likelihood MCMC results reported here. Those results are presented as an independent cross-check, alongside the diagnostic-only CMB distance-prior view already used elsewhere.
\end{mdframed}


\section{Cosmological validation: CMB, BAO, and local H0}
\label{sec:cosmo}
This section documents a matched, reproducible cosmological comparison between $\Lambda$CDM and MTDF using Planck 2018 TTTEEE likelihoods and DESI BAO, followed by an explicit test of the SH0ES local H0 prior under an MTDF void-enhancement mapping. All quantities below are reported exactly as evaluated by the likelihood code used for the MCMC runs.

\begin{mdframed}[linecolor=accent,linewidth=0.5pt,backgroundcolor=callout-bg,leftline=true,rightline=true,topline=true,bottomline=true,innerleftmargin=8pt,innerrightmargin=8pt,innertopmargin=6pt,innerbottommargin=6pt]
Methodological note: this section reports a matched CMB+BAO comparison (Planck TTTEEE plus DESI BAO) and then adds a local-H0 interpretation layer (SH0ES). The SH0ES layer is treated as a test of whether one parameter set can jointly accommodate global and local expansion in a controlled way. These results are reported here and are not included in the dashboard strict scorecard. The definitive full Planck plik TTTEEE + low-$\ell$ + lensing hard-falsifier is reported separately in Section 2.4.
\end{mdframed}

\subsection{Matched CMB + BAO comparison (Planck 2018 TTTEEE + DESI BAO)}
\textbf{Implementation-era note.} The summary runs in Sections 2.1 and 2.2 predate the v2 growth-sector correction and are retained as historical bookkeeping context; their distance-sector conclusions are insensitive to the corrected growth source, and any growth-sensitive reading defers to the sealed v2 results of Section 2.4. They are regenerated against the corrected implementation at the derived-products step.

We ran joint MCMC fits for $\Lambda$CDM and MTDF on identical data and nuisance settings. MTDF adds one additional parameter (k\_\allowbreak{}f) while keeping the standard six cosmological parameters. In this CMB+BAO summary run, MTDF matches the fit quality to within $\Delta$$\chi$$^{2}$ $\approx$ 1.25. This number refers only to the CMB+BAO summary table below; the full Planck plik TTTEEE + low-$\ell$ + lensing comparison is reported in Section 2.4 (sealed v2 chains); the completed comparison against the matched $\Lambda$CDM control is given in Section 2.4.4.

\begin{table}[H]
\centering
\footnotesize
\begin{tabularx}{\linewidth}{l>{\RaggedRight\arraybackslash}X>{\RaggedRight\arraybackslash}X>{\RaggedRight\arraybackslash}X>{\RaggedRight\arraybackslash}X>{\RaggedRight\arraybackslash}X}
\toprule
\textbf{Model} & \textbf{$\chi$$^{2}$ Total} & \textbf{$\chi$$^{2}$ Planck} & \textbf{$\chi$$^{2}$ BAO} & \textbf{H0 (km/\allowbreak{}s/\allowbreak{}Mpc)} & \textbf{k\_\allowbreak{}f} \tabularnewline
\midrule
$\Lambda$CDM & 1018.57 & 1004.78 & 13.79 & 68.03 (+0.43/\allowbreak{}-0.29) & n/\allowbreak{}a \tabularnewline
MTDF & 1019.82 & 1005.23 & 14.59 & 68.13 (+0.54/\allowbreak{}-0.35) & 1.00 (+0.09/\allowbreak{}-0.05) \tabularnewline
\bottomrule
\end{tabularx}
\end{table}

Interpretation: this indicates that, within the stated implementation, introducing k\_\allowbreak{}f does not produce a statistically meaningful degradation of the global CMB+BAO fit. Any claimed advantage must then come from additional, independently testable phenomena, such as the local-environment mapping tested in Part I.

\subsection{SH0ES prior and the global vs local H0 test}
We then evaluated the SH0ES H0 prior as an additional constraint. A naive application (treating SH0ES as directly measuring the global H0) does not relieve the tension. In contrast, the MTDF two-component interpretation treats Planck as probing a global effective expansion rate and SH0ES as probing a locally void-enhanced rate. Under that mapping, the residual tension is reduced substantially in the current run.

\begin{table}[H]
\centering
\footnotesize
\begin{tabularx}{\linewidth}{l>{\RaggedRight\arraybackslash}X>{\RaggedRight\arraybackslash}X>{\RaggedRight\arraybackslash}X}
\toprule
\textbf{Model} & \textbf{$\chi$$^{2}$\_\allowbreak{}SH0ES} & \textbf{$\chi$$^{2}$ Total} & \textbf{Residual tension} \tabularnewline
\midrule
$\Lambda$CDM + SH0ES & 14.2 & 1034.7 & 3.8$\sigma$ \tabularnewline
MTDF (naive) + SH0ES & 14.0 & 1033.4 & 3.8$\sigma$ \tabularnewline
MTDF (with void mapping) + SH0ES & 1.64 & 1019.4 & 1.3$\sigma$ \tabularnewline
\bottomrule
\end{tabularx}
\end{table}

\textbf{Conditional H0 chain (Phase 5 bookkeeping):}

\begin{itemize}
  \item Global H0 (sealed v2 full Planck posterior, Section 2.4.2): 67.44 km/\allowbreak{}s/\allowbreak{}Mpc
  \item Void enhancement factor: S0 = 1.084 (the algebraic identity S0 = $\alpha$$\sqrt{\Omega_{\rm DE}}$, evaluated at KBC-level void depth)
  \item Implied local value: H\_\allowbreak{}local $\approx$ 73.1 km/\allowbreak{}s/\allowbreak{}Mpc, conditional on a deep KBC-like local underdensity; shallower void reconstructions give roughly 70 to 71 km/\allowbreak{}s/\allowbreak{}Mpc
  \item SH0ES observed: 73.04 $\pm$ 1.04 km/\allowbreak{}s/\allowbreak{}Mpc
\end{itemize}

MTDF does not predict the local void itself: the depth of the local underdensity is an observational input, and the implied local value is conditional on it.

Provenance of the void coupling: the parameter $\alpha$ = 1.30 entering the enhancement identity S0 = $\alpha$$\sqrt{\Omega_{\rm DE}}$ is a frozen, void-motivated phenomenological coupling, not an independently measured quantity. The void-depth inputs that motivate it come from documented reconstructions of the local volume (the Tully et al. Cosmicflows-3 velocity field, indicating an outflow of roughly +13.1\%, and Boehringer-Jasche-Lavaux-type density reconstructions indicating a 10 to 20\% local underdensity), and these reconstructions are not fully independent of the distance ladder itself, so a circularity caveat applies when the mapping is compared against SH0ES. The S0 to H\_\allowbreak{}local mapping is therefore a consistency statement, not a prediction: given a deep KBC-like local void, the global posterior implies a local value near 73.1 km/\allowbreak{}s/\allowbreak{}Mpc, while shallower reconstructions imply roughly 70 to 71 km/\allowbreak{}s/\allowbreak{}Mpc.

Within MTDF, the Hubble tension is not treated as a measurement error or an unmodelled systematic; it is interpreted as an environmental split between a global expansion rate and a locally inferred, environment-enhanced rate. In the implementation tested here, CMB plus BAO constrain a global-average value near 67 to 68 km/\allowbreak{}s/\allowbreak{}Mpc, while the local distance ladder probes a void-weighted value near 73 to 74 km/\allowbreak{}s/\allowbreak{}Mpc. Under these assumptions, both measurements can be accommodated as probes of different effective regimes within the same field model.

\begin{itemize}
  \item Planck (CMB) probes the global-average H0, approximately 67 to 68 km/\allowbreak{}s/\allowbreak{}Mpc.
  \item SH0ES (local ladder) probes a locally weighted H0 in low-density environments, approximately 73 to 74 km/\allowbreak{}s/\allowbreak{}Mpc.
  \item Interpretation: both are consistent if they weight different environmental regimes.
\end{itemize}

This is falsifiable: the void-enhancement mapping must remain consistent under future void catalogues, alternative low-z calibrators, and expanded local samples. If the mapping fails, the interpretation is constrained without affecting the matched CMB+BAO fit result in Section 2.1.

\subsection{Conservative fit philosophy}
To avoid artefactual improvements, we report likelihood contributions transparently and keep optional layers separated. In particular, the SH0ES reduction above is not obtained by tuning the global CMB+BAO likelihood, but by adding a separately testable environmental mapping whose sign and scale are constrained by the Part I supernova void analysis.

The growth behaviour inferred from the sealed v2 full Planck likelihood (an enhanced-growth posture, with $\sigma$\textsubscript{8} about 1.9$\sigma$ above the $\Lambda$CDM baseline in posterior units) motivates the registered growth-history and forward-modelling programmes (GH-1, FWD-S8) beyond the scope of this work.

\subsection{Full Planck plik TTTEEE + lensing validation (v2 chains, sealed)}
\label{sec:phase5}

\begin{mdframed}[linecolor=accent,linewidth=0.5pt,backgroundcolor=callout-bg,leftline=true,rightline=true,topline=true,bottomline=true,innerleftmargin=8pt,innerrightmargin=8pt,innertopmargin=6pt,innerbottommargin=6pt]
\textbf{Scope of the cosmological test (mandatory reading for every citation of these results).} The v2 run is a full Planck 2018 likelihood test of the \emph{corrected MTDF effective cosmological implementation}: the corrected covariant $\mu$(a)/EFE sector, with the high-Q carrier represented by its CDM-like limit (standard pressureless CDM as proxy). It is not a Boltzmann implementation of all nine degrees of freedom of the frozen action; the full-action CMB perturbation implementation is a named future calculation. Success at this level discharges the CMB confrontation at the effective level only.
\end{mdframed}

\begin{mdframed}[linecolor=accent,linewidth=0.5pt,backgroundcolor=callout-bg,leftline=true,rightline=true,topline=true,bottomline=true,innerleftmargin=8pt,innerrightmargin=8pt,innertopmargin=6pt,innerbottommargin=6pt]
\textbf{Supersession notice (v1 chains retired).} The Phase-5 chain report published in V74 (posterior table with $\sigma$\textsubscript{8} = 0.790, S\textsubscript{8} = 0.802, $\Delta$$\chi$$^{2}$ = +0.63, k\textsubscript{f} = 0.50 $\pm$ 0.36, and the associated correlation and quartile tables) was produced with the v1 growth implementation, which carried a documented gauge artifact: the growth modification was applied to the synchronous-gauge 00-source, breaking a near-cancellation and flipping the sign of the growth response. Those numbers are void and are retained only in the archived V74 release as a historical record. Everything below is computed from the corrected v2 implementation (covariant $\mu$ in newtonian gauge; frozen source \texttt{{\small\protect\path{perturbations.c}}}, hash-pinned in the flight package).
\end{mdframed}

\subsubsection{Run provenance and convergence}
The v2 production run used the frozen flight configuration: full Planck 2018 plik TTTEEE + low-$\ell$ TT/EE + lensing, the corrected class\_\allowbreak{}mtdf binary, four chains, identical priors and stop rules to the frozen v2 YAML set, with k\textsubscript{f} the only free MTDF parameter ($\alpha$, $\beta$\textsubscript{eos}, $\tau$ frozen). The run reached the frozen hard-cap stop rule of 200,000 accepted steps per chain (800,000 total) after \~{}66 h, with R$-$1 = 0.0085 (target 0.01, satisfied) and R$-$1\textsubscript{CL} = 0.060 (strict bound-level target 0.02 not reached: 95 percent bounds carry extra sampling noise; means and peaks are solid). All interruptions auto-resumed from checkpoint with no sample loss.

\subsubsection{Sealed posterior comparison (k\textsubscript{f} free)}
\begin{table}[H]
\centering
\footnotesize
\begin{tabularx}{\linewidth}{l>{\RaggedRight\arraybackslash}X>{\RaggedRight\arraybackslash}X>{\RaggedRight\arraybackslash}X}
\toprule
\textbf{Parameter} & \textbf{$\Lambda$CDM baseline} & \textbf{MTDF v2 (mean $\pm$ $\sigma$)} & \textbf{Shift} \tabularnewline
\midrule
$\sigma$\textsubscript{8} & 0.8101 & 0.82651 $\pm$ 0.00623 & +0.0164 (+1.91$\sigma$, posterior units) \tabularnewline
H\textsubscript{0} (km/\allowbreak{}s/\allowbreak{}Mpc) & 67.382 & 67.444 $\pm$ 0.546 & +0.06 (CMB-side H\textsubscript{0} unchanged) \tabularnewline
$\Omega$\textsubscript{m} & -- & 0.31533 $\pm$ 0.00771 & negligible \tabularnewline
n\textsubscript{s} & -- & 0.96426 & negligible \tabularnewline
ln(10\textsuperscript{10}A\textsubscript{s}) & -- & 3.03962 & negligible \tabularnewline
k\textsubscript{f} & -- & 0.685 $\pm$ 0.500 & consistent with 0; theory value 1 at $-$0.63$\sigma$ \tabularnewline
\bottomrule
\end{tabularx}
\end{table}

Sealed v2 reporting set. The $\Lambda$CDM baseline values are the same-pipeline v2-era baselines used by the sealed adjudication; the full baseline column and all dataset-resolved comparisons are regenerated from the machine-readable v2 results package together with the matched $\Lambda$CDM control chain (see 2.4.4).

\textbf{Posture (sealed, pre-registered three-posture rule): S8-WORSENED.} The corrected covariant $\mu$(a) enhances late growth by about +2 percent, and the full-likelihood chain confirms this end to end: $\sigma$\textsubscript{8} lands about 2$\sigma$ (posterior units) above the $\Lambda$CDM baseline. The corrected theory therefore mildly worsens the comparison with published weak-lensing clustering amplitudes relative to $\Lambda$CDM. In native units the growth-sector posterior gives S\textsubscript{8} = 0.8473 $\pm$ 0.0138 (free) and 0.8501 $\pm$ 0.0132 (k\textsubscript{f} = 1). The importance-reweighted k\textsubscript{f} = 1 rider gives the same posture ($\sigma$\textsubscript{8} shift +0.0173, +2.07$\sigma$), with f\textsubscript{kick} = 0.0033 and $\mu$(z\textsubscript{rec}) = 1.00002 (CMB-inert), as derived. Standard-parameter shifts are tiny and the CMB-side H\textsubscript{0} is unchanged, consistent with the retired ``H\textsubscript{0} derived'' claim remaining retired. Not affected: the galactic-sector validations (RAR/BTFR), c\textsubscript{T} = 1, $\eta$ = 1, the frozen 9-degree-of-freedom action, and the finding that Planck alone cannot pin the EFE amplitude.

\subsubsection{Coupling posterior}
The k\textsubscript{f} posterior is broad, k\textsubscript{f} = 0.685 $\pm$ 0.500: both k\textsubscript{f} = 0 ($\Lambda$CDM) and k\textsubscript{f} = 1 (full coupling) are compatible, so Planck establishes consistency, not support, exactly as in the superseded v1 chains (whose k\textsubscript{f}-unconstrained finding survives). Discrimination requires low-redshift probes (Part I environment signal; growth history GH-1).

\subsubsection{Model comparison against the matched $\Lambda$CDM control}
The matched $\Lambda$CDM control chain shares the entire v2 infrastructure (binary, likelihood set, priors on all shared parameters, four-chain profile, stop rules); its configuration differs from the frozen v2 YAML in exactly three deltas that remove the MTDF theory keys. Its posterior reproduces the reference Planck-$\Lambda$CDM universe at full-posterior level: $\sigma$\textsubscript{8} = 0.81036 $\pm$ 0.00598, per-sample S\textsubscript{8} = 0.83004 $\pm$ 0.01278, H\textsubscript{0} = 67.389 $\pm$ 0.544, with all shared-parameter shifts against the reference baselines at the 10\textsuperscript{$-$4} level. This validates the corrected binary's $\Lambda$CDM limit end to end and qualifies the control as the injection-recovery truth for the preregistered forward-modelling programme.

\textbf{Matched $\Delta$$\chi$$^{2}$.} At the bounded best fits (minimisation seeded from the best sampled point of each chain; both converged), MTDF gives $-$log L = 1385.78 ($\chi$$^{2}$ = 2771.56, at k\textsubscript{f} = 0.476) against $\Lambda$CDM $-$log L = 1386.42 ($\chi$$^{2}$ = 2772.84):

\begin{center}
\textbf{$\Delta$$\chi$$^{2}$ (MTDF $-$ $\Lambda$CDM) = $-$1.29}
\end{center}

\emph{This is a maximum-likelihood comparison, not a Bayes factor or Bayesian evidence calculation.} MTDF carries one additional free parameter (k\textsubscript{f}), and a gain of one to two $\chi$$^{2}$ points per extra degree of freedom is the no-information expectation for nested fits; the Akaike-corrected difference is $\Delta$AIC = +0.71, favouring neither model at any meaningful level. (A BIC-style penalty depends on an adopted effective Planck data count and is not quoted; under any conventional count the verdict is unchanged.) The full Planck 2018 likelihood therefore does not distinguish the corrected MTDF from $\Lambda$CDM. The bounded MTDF best fit sits at k\textsubscript{f} = 0.476, independently consistent with the broad k\textsubscript{f} posterior of 2.4.3.

\begin{table}[H]
\centering
\footnotesize
\begin{tabularx}{\linewidth}{l>{\RaggedRight\arraybackslash}X>{\RaggedRight\arraybackslash}X>{\RaggedRight\arraybackslash}X}
\toprule
\textbf{Dataset} & \textbf{$\Lambda$CDM $\chi$$^{2}$} & \textbf{MTDF $\chi$$^{2}$} & \textbf{$\Delta$} \tabularnewline
\midrule
Planck 2018 lowl TT & 23.77 & 22.29 & $-$1.48 \tabularnewline
Planck 2018 lowl EE & 396.27 & 396.47 & +0.20 \tabularnewline
Planck 2018 plik TTTEEE & 2344.00 & 2343.80 & $-$0.20 \tabularnewline
Planck 2018 lensing & 8.80 & 8.99 & +0.19 \tabularnewline
\textbf{Total} & \textbf{2772.84} & \textbf{2771.56} & \textbf{$-$1.29} \tabularnewline
\bottomrule
\end{tabularx}
\end{table}

Dataset-resolved $\chi$$^{2}$ (= 2 $\times$ $-$log L per component) at the bounded best fits. Every component agrees to about one point; the small MTDF gain sits in lowl TT, the late-ISW region, consistent with a late-time-only physics difference. Per-sample growth-sector amplitudes: MTDF S\textsubscript{8} = 0.8473 $\pm$ 0.0138 (k\textsubscript{f} free) and 0.8501 $\pm$ 0.0132 (k\textsubscript{f} = 1) against the control's 0.83004 $\pm$ 0.01278.

Convention note: sampler-reported chain minima are minus-log-posterior values, which include each model's constant flat-prior volume term and are not comparable across models of different dimensionality; all comparisons above are likelihood-only, taken from the bounded minima and the per-dataset $\chi$$^{2}$ components.

\subsubsection{Detailed component analyses}
The correlation, quartile, and component-resolved analyses of the superseded v1 report are regenerated from the v2 chains as part of the machine-readable results package that accompanies the sealed control-chain comparison; per the single-numerical-source rule, this paper does not carry hand-copied component tables ahead of that package.

\begin{mdframed}[linecolor=accent,linewidth=0.5pt,backgroundcolor=callout-bg,leftline=true,rightline=true,topline=true,bottomline=true,innerleftmargin=8pt,innerrightmargin=8pt,innertopmargin=6pt,innerbottommargin=6pt]
\textbf{Reproducibility.} The v2 output tarball is hash-pinned (sha256 95ce9385\ldots{}), produced by a frozen, sha256-pinned run configuration and binary, adjudicated by the sealed analysis script (\texttt{{\small\protect\path{analyze_v2.py}}}) under a pre-registered three-posture rule with the adjudication executed only on the completed chain. The v1 Phase-5 artifacts remain in the archived V74 release as historical record.
\end{mdframed}

\section{High-energy phenomenology}
\label{sec:high-energy}
These subsections are exploratory hypothesis statements, included solely to define potential falsifiers. They carry no evidential weight in the MTDF programme, and no result elsewhere in the suite depends on them.

\subsection{Black holes as stress singularities}
\begin{mdframed}[linecolor=accent,linewidth=0.5pt,backgroundcolor=callout-bg,leftline=true,rightline=true,topline=true,bottomline=true,innerleftmargin=8pt,innerrightmargin=8pt,innertopmargin=6pt,innerbottommargin=6pt]
Hypothesis (to be tested).
\end{mdframed}

\subsubsection{Observation}
Compact objects and accretion systems exhibit jets, variability, and energetic outflows that are not always easy to connect to a single universal mechanism across mass scales.

\subsubsection{MTDF mechanism}
In this hypothesis, extreme curvature environments correspond to extreme stress concentrations. Discharge, relaxation, or topological transitions could manifest as energetic events or persistent jets when fed by accretion driven stress loading.

$$
\frac{d}{dt}\,u_{\mathrm{tensor}} \approx \mathcal{P}_{\mathrm{load}} - \mathcal{P}_{\mathrm{relax}}
$$

\subsubsection{Predictions}
\begin{itemize}
  \item Event energetics and recurrence times should correlate with an inferred stress loading rate proxy, not only with mass.
  \item Systems in similar accretion states but different environments could show systematic differences if background stress matters.
  \item If discharge is real, there should be identifiable precursors in variability statistics.
\end{itemize}

\subsubsection{How to test}
\begin{itemize}
  \item Search for precursor signatures in power spectra and time series of well monitored AGN and X ray binaries.
  \item Test whether jet power correlates better with stress proxies than with halo or environment measures used in standard models.
  \item Use population studies to look for bimodality that could indicate threshold discharge behaviour.
\end{itemize}

\begin{mdframed}[linecolor=accent,linewidth=0.5pt,backgroundcolor=callout-bg,leftline=true,rightline=true,topline=true,bottomline=true,innerleftmargin=8pt,innerrightmargin=8pt,innertopmargin=6pt,innerbottommargin=6pt]
Methodological note: Even if the hypothesis fails, the exercise clarifies which high energy observations are most diagnostic for stress based extensions.
\end{mdframed}

\subsection{Fast radio bursts as localised stress oscillations}
\begin{mdframed}[linecolor=accent,linewidth=0.5pt,backgroundcolor=callout-bg,leftline=true,rightline=true,topline=true,bottomline=true,innerleftmargin=8pt,innerrightmargin=8pt,innertopmargin=6pt,innerbottommargin=6pt]
Hypothesis (to be tested).
\end{mdframed}

\subsubsection{Observation}
FRBs show rapid time structure, diverse repetition behaviour, and complex environments. Many models exist. A stress based hypothesis must compete with strong alternatives.

\subsubsection{MTDF mechanism}
Localised stress oscillations in a magnetised plasma environment could act as a trigger or modulator, producing bursts when a threshold is crossed. The stress term is not the full mechanism; it is a proposed enabling channel.

\subsubsection{Predictions}
\begin{itemize}
  \item If stress oscillations contribute, repetition rates may correlate with large scale environment or with local structure proxies.
  \item There could be a characteristic distribution of waiting times consistent with relaxation and recharge dynamics.
  \item Bursts might show correlated behaviour with other tracers of rapid structural change in the host region.
\end{itemize}

\subsubsection{How to test}
\begin{itemize}
  \item Use well localised FRB samples to test for environment correlations beyond host galaxy properties.
  \item Fit waiting time distributions with recharge style models and compare across repeating and non repeating classes.
  \item Search for cross correlations with transient high energy or optical activity in hosts where data exist.
\end{itemize}

\section{Lensing and merging clusters}
\label{sec:lensing}
Lensing is a high leverage domain because it probes the gravitational field without relying on tracer dynamics. Any extension that aims to replace dark matter must ultimately be compatible with lensing on galaxy and cluster scales.

\subsection{Lensing consistency as a gatekeeper test}
\begin{mdframed}[linecolor=accent,linewidth=0.5pt,backgroundcolor=callout-bg,leftline=true,rightline=true,topline=true,bottomline=true,innerleftmargin=8pt,innerrightmargin=8pt,innertopmargin=6pt,innerbottommargin=6pt]
Cross-cutting requirement: any field-based alternative to dark matter must reproduce the empirical weak-lensing mass maps, including in merging systems. This section states the requirement and the associated falsifiers for MTDF-based extensions.
\end{mdframed}

\subsubsection{Observation}
Weak lensing and strong lensing provide mass map constraints that are often described as evidence for collisionless dark matter when compared to baryonic tracers.

\subsubsection{MTDF mechanism}
A stress field contribution must generate lensing potentials consistent with observed shear and deflection statistics, using the same parameter set used for other scales, and without hiding tuning freedom.

\subsubsection{Predictions}
\begin{itemize}
  \item Cluster scale lensing maps should be reproducible when baryon distributions and stress field response are modelled transparently.
  \item Merging cluster morphologies are a particularly sharp test, because they separate baryonic gas and inferred lensing mass.
  \item If the stress field has memory, transient lensing signatures might exist in unrelaxed systems.
\end{itemize}

\subsubsection{How to test}
\begin{itemize}
  \item Define a benchmark set of clusters and merging systems and specify in advance which modelling degrees of freedom are permitted.
  \item Run blinded analyses comparing stress field based reconstructions against standard mass models.
  \item Require that the same parameter set fits both galaxy kinematics and lensing within declared uncertainties.
\end{itemize}

\subsection{Cluster scale mass bias and stress supported hydrostatics}
\begin{mdframed}[linecolor=accent,linewidth=0.5pt,backgroundcolor=callout-bg,leftline=true,rightline=true,topline=true,bottomline=true,innerleftmargin=8pt,innerrightmargin=8pt,innertopmargin=6pt,innerbottommargin=6pt]
Note: Extension. This section reframes known cluster systematics as potential stress support signatures. It is not asserted as a completed fit.
\end{mdframed}

\subsubsection{Observation}
Galaxy clusters are routinely weighed by several methods that do not always agree. X-ray and SZ based hydrostatic masses can differ from weak-lensing masses,
and the mismatch is often discussed as "hydrostatic mass bias" driven by non-thermal pressure support, turbulence, and departures from equilibrium.

\subsubsection{MTDF mechanism}
In MTDF, the stress field can contribute an additional, non-baryonic support term that affects both the effective force balance of the intracluster medium and
the lensing potential inferred from light deflection. In this framing, part of the apparent mass discrepancy is not missing matter but stress support that is not
captured by purely thermal hydrostatic modelling.

\subsubsection{Predictions}
\begin{itemize}
  \item \textbf{Dynamical state dependence:} the mass mismatch should be larger in disturbed or recently merged systems where stress gradients are expected to be stronger.
  \item \textbf{Spatial signatures:} residuals should have a characteristic radial profile linked to the finite-range stress response and relaxation, rather than the profile expected from an additional collisionless component.
  \item \textbf{Cross-probe consistency:} a single stress-support prescription should improve consistency between X-ray, SZ, galaxy dynamics, and weak lensing without retuning the fundamental parameter set.
\end{itemize}

\subsubsection{How to test}
\begin{itemize}
  \item Use cluster samples with both weak lensing maps and high quality X-ray or SZ pressure profiles. Quantify the residual force balance required beyond thermal support and compare it across dynamical states.
  \item Test whether the inferred "extra support" correlates with large scale environment and merger indicators (substructure measures, centroid shifts, shock features).
  \item State falsifiers explicitly: if residuals are fully explained by conventional non-thermal pressure models with no additional coherent term, then this stress-support extension is constrained to be negligible at cluster scales.
\end{itemize}

\subsection{Merging clusters and mass map offsets as a critical test}
\begin{mdframed}[linecolor=accent,linewidth=0.5pt,backgroundcolor=callout-bg,leftline=true,rightline=true,topline=true,bottomline=true,innerleftmargin=8pt,innerrightmargin=8pt,innertopmargin=6pt,innerbottommargin=6pt]
High-priority falsifier: merging-cluster lensing offsets are a decisive test for any alternative to dark matter. This subsection states the quantitative target and what would rule out the proposed stress-peak interpretation.
\end{mdframed}

\subsubsection{Observation}
Some merging clusters exhibit weak-lensing mass peaks that are offset from the hot gas traced by X-ray emission. In the standard interpretation, this is evidence for
a dominant collisionless component that separates from baryons during the merger.

\subsubsection{MTDF mechanism}
A viable MTDF extension must allow transient, spatially localised stress concentrations during mergers. The candidate response invokes the finite relaxation time $\tau$ = 13 Gyr: a stress configuration sourced over Gyr timescales need not track the gas through a \~{}100 Myr merger, and may remain associated with the pre-merger mass distribution traced by the collisionless galaxies, producing an effective lensing contribution that does not coincide with the gas distribution at a given snapshot. This mechanism is stated here as a hypothesis; no simulation supports it yet, and a quantitative test requires time-dependent field evolution in an N-body setting. It is the same open falsifier described in the gravity-sector companion (MTDF\_\allowbreak{}03 Section 6.5).

\subsubsection{Predictions}
\begin{itemize}
  \item \textbf{Time dependence:} offsets should depend on merger phase. The relative position of lensing peaks, galaxies, and gas should evolve predictably rather than remain static.
  \item \textbf{Pattern consistency across systems:} the distribution of offsets across many mergers should follow a stress-propagation pattern tied to shock geometry and relaxation, not an arbitrary set of displacements.
  \item \textbf{Multi-wavelength linkage:} stress peak locations should correlate with merger features such as shocks and shear flows that indicate rapid stress loading.
\end{itemize}

\subsubsection{How to test}
\begin{itemize}
  \item Assemble a catalogue of merging clusters with comparable lensing reconstructions and X-ray maps. Compare offset statistics against merger stage indicators.
  \item Use forward simulations that include a stress field component to test whether observed offset amplitudes and morphologies can be reproduced without introducing a new particle component.
  \item \textbf{Hard falsifier:} if realistic stress-field simulations cannot produce the observed offset morphologies and amplitudes while preserving the successful galaxy-scale behaviour, the MTDF extension fails at cluster scale.
\end{itemize}

\section{Structured test programme}
\label{sec:tests}
To keep this companion scientifically useful, each hypothesis should be paired with a small set of near term tests and a longer term roadmap. A suggested ordering:

\textbf{Additional high leverage items for this part:}

\begin{itemize}
  \item Cluster mass bias: joint weak-lensing plus X-ray or SZ analyses with explicit residual force-balance accounting.
  \item Merging clusters: offset statistics across a controlled sample, with phase indicators and forward modelling.
  \item Cross-check: require that any cluster-scale extension does not spoil the galaxy-scale relations when the same fundamental parameters are held fixed.
\end{itemize}

\textbf{Link:} Non-merging cluster-scale forward predictions and baryon completion checks are reported in the gravity-sector companion paper (see Step 21 in \textit{MTDF\_\allowbreak{}03 (Companion Part II)}~[see companion paper]). The merger offset morphology programme remains an open, higher-complexity falsifier.

\begin{enumerate}
  \item Start with high leverage, low ambiguity probes: lensing consistency, isotropy constraints, and standard ruler cross checks.
  \item Move to environment split analyses on well characterised datasets.
  \item Only then elevate high energy phenomenology claims, after basic constraints are satisfied.
\end{enumerate}

\begin{mdframed}[linecolor=accent,linewidth=0.5pt,backgroundcolor=callout-bg,leftline=true,rightline=true,topline=true,bottomline=true,innerleftmargin=8pt,innerrightmargin=8pt,innertopmargin=6pt,innerbottommargin=6pt]
\textbf{Reproducibility and interactive materials.}\begin{itemize}
  \item Validation suite appendix: \textit{{\small\protect\path{MTDF_06_Validation_Suite_Appendix.html}}}~[see companion paper]
  \item Interactive dashboard: {\small\protect\path{Validation_Dashboard_V75.html}}
\end{itemize}
\end{mdframed}

\section{Data and code availability}
The complete MTDF reproducibility package, including the modified CLASS solver
(\texttt{class\_\allowbreak{}mtdf}), the V19 strict validation workbook, the
\texttt{{\small\protect\path{run_validate.py}}} pipeline, the gravity-sector artefact manifest,
and the Phase 2 CosmoPower emulator cache, is archived on Zenodo
(concept DOI \href{https://doi.org/10.5281/zenodo.19741058}{{\small\protect\path{10.5281/zenodo.19741058}}};
this release, \texttt{v1.1.7}, DOI
\href{https://doi.org/10.5281/zenodo.21525855}{{\small\protect\path{10.5281/zenodo.21525855}}})
and mirrored on GitHub at
\href{https://github.com/IMesche/MTDF}{github.com/\allowbreak{}IMesche/\allowbreak{}MTDF}.
The \texttt{v1.1.7} Git tag corresponds to the archived Zenodo release. All
numerical results reported here can be regenerated from the shipped workbook
and scripts following the procedure in the top-level \texttt{{\small\protect\path{README.md}}}.

\section{References}
{\scriptsize
\begin{thebibliography}{11}
\bibitem{ref1} Aghanim, N. et al. (Planck Collaboration) 2020, "Planck 2018 results. VI. Cosmological parameters", \emph{A\&A}, 641, A6. DOI: \href{https://doi.org/10.1051/0004-6361/201833910}{{\small\protect\path{10.1051/0004-6361/201833910}}}
\bibitem{ref2} Aghanim, N. et al. (Planck Collaboration) 2020, "Planck 2018 results. VIII. Gravitational lensing", \emph{A\&A}, 641, A8. DOI: \href{https://doi.org/10.1051/0004-6361/201833886}{{\small\protect\path{10.1051/0004-6361/201833886}}}
\bibitem{ref3} DESI Collaboration 2025, "DESI 2024 VI: Cosmological Constraints from BAO", \emph{JCAP}, 2025, 021. DOI: \href{https://doi.org/10.1088/1475-7516/2025/02/021}{{\small\protect\path{10.1088/1475-7516/2025/02/021}}}
\bibitem{ref4} Riess, A. G. et al. 2022, "A Comprehensive Measurement of the Local Value of the Hubble Constant", \emph{ApJL}, 934, L7. DOI: \href{https://doi.org/10.3847/2041-8213/ac5c5b}{{\small\protect\path{10.3847/2041-8213/ac5c5b}}}
\bibitem{ref5} Asgari, M. et al. (KiDS Collaboration) 2021, "KiDS-1000 Cosmology: Cosmic shear constraints", \emph{A\&A}, 645, A104. DOI: \href{https://doi.org/10.1051/0004-6361/202039070}{{\small\protect\path{10.1051/0004-6361/202039070}}}
\bibitem{ref6} Amon, A. et al. (DES Collaboration) 2022, "Dark Energy Survey Year 3 results: Cosmology from cosmic shear", \emph{Phys. Rev. D}, 105, 023514. DOI: \href{https://doi.org/10.1103/PhysRevD.105.023514}{{\small\protect\path{10.1103/PhysRevD.105.023514}}}
\bibitem{ref7} Torrado, J. \& Lewis, A. 2021, "Cobaya: Code for Bayesian Analysis of hierarchical physical models", \emph{JCAP}, 2021, 057. DOI: \href{https://doi.org/10.1088/1475-7516/2021/05/057}{{\small\protect\path{10.1088/1475-7516/2021/05/057}}}
\bibitem{ref8} Clowe, D. et al. 2006, "A Direct Empirical Proof of the Existence of Dark Matter", \emph{ApJL}, 648, L109. DOI: \href{https://doi.org/10.1086/508162}{{\small\protect\path{10.1086/508162}}}
\bibitem{ref9} Ade, P. A. R. et al. (Planck Collaboration) 2016, "Planck 2015 results. XXIV. Cosmology from SZ cluster counts", \emph{A\&A}, 594, A24. DOI: \href{https://doi.org/10.1051/0004-6361/201525833}{{\small\protect\path{10.1051/0004-6361/201525833}}}
\bibitem{ref10} Alam, S. et al. (eBOSS Collaboration) 2021, "Completed SDSS-IV eBOSS: Cosmological implications", \emph{Phys. Rev. D}, 103, 083533. DOI: \href{https://doi.org/10.1103/PhysRevD.103.083533}{{\small\protect\path{10.1103/PhysRevD.103.083533}}}
\bibitem{ref11} B\"ohringer, H., Chon, G. \& Collins, C. A. 2020, "Observational evidence for a local underdensity in the Universe and its effect on the measurement of the Hubble constant", \emph{A\&A}, 633, A19. DOI: \href{https://doi.org/10.1051/0004-6361/201936400}{{\small\protect\path{10.1051/0004-6361/201936400}}}
\end{thebibliography}
}


\end{document}
